%% kapa: exam=2 round=34 year=2023 period=1 subject=appraisal_practice subject_ko=감정평가실무 seq=3 number=3 points=20
%% source: Q-Net 감정평가사 기출문제 (cst00309) artlSeq=5212168 file=raw/1008/5212168/1교시 23년도 감정평가사 2차시험_감정평가실무.pdf
%% conversion: kordoc 4.24.51 -> tools/gen2.py (auto); review=proofread-gpt-sol; flags=none
\begin{longquestion}{【문제 3】}
감정평가사 甲은 수도권 북동부 소재 A택지개발지구 내에 분양중인 상업용지 개발과 관련된 자문을 의뢰받았다. 제시된 자료를 활용하여 다음 물음에 답하시오. (20점) \par
\medskip\noindent 물음 1) 해당 토지의 개발이 가능한 임대료(Feasibility Rent)를 구하고, 그 결과를 바탕으로 현재의 개발 타당성을 판단하시오. (10점) \par
\medskip\noindent 물음 2) 현재 개발이 타당하지 않다면, 개발사업에 소요되는 기간을 고려할 경우 얼마 후 사업에 착수하는 것이 타당한지 그 시기를 구하시오. (10점) \par
< 자료 1 > 기본적 사항 \par
1. 토지면적: 3,000m² \par
2. 용적률: 250\% \par
3. 토지 분양가격(부대비용 포함): @6,000,000원/m² \par
4. 건축비(부대비용 포함): @1,800,000원/m² \par
5. 건물 전용률: 60\% \par
< 자료 2 > 시장조사 내용 요약 \par
1. 시장임대료(전유면적 기준): 보증금 @400,000원/m², 월임료 @40,000원/m²(전환율: 연 9.0\%) \par
2. 운영경비 \par
1) 연간 조세공과(보유세 등): 토지가격 및 건물가격의 0.25\% \par
2) 연간 보험료: 건물가격의 0.2\% \par
3) 변동비용: 유효총수익(EGI)의 5.0\% \par
4) 건물설비 중 엘리베이터(현재 재조달원가 1.2억원)는 내용연수 12년으로 내용연수 만료 시에 신품으로 대체하여야 함 \par
3. 공실률: 6.0\% \par
4. 할인율: 연 7.0\% \par
< 자료 3 > 시장전망 요약 \par
1. 최근 금리 상승 등에 따른 영향으로 향후 토지 분양가격 및 건축비, 인근지역의 상업용 부동산가격은 보합세를 보일 것으로 예상됨 \par
2. 개발 착수 후 사용승인까지는 9개월이 소요될 것으로 예상되며, 개발기간 중의 제반비용(금융비용 등)은 고려하지 않음 \par
3. 향후 주변 신축 아파트의 지속적인 입주로 상업시설공간에 대한 안정적인 수요가 있어 사용승인과 더불어 현재 공실률 수준의 임대가 가능할 것으로 예상되며, 시장임대료는 연 5.0\%, 운영경비(엘리베이터 가격 포함)는 연 2.0\% 상승할 것으로 예상됨 \par
\end{longquestion}
